\documentclass{amsart}
% For dealing with references we use the comment environment
\usepackage{verbatim}
\newenvironment{reference}{\comment}{\endcomment}
%\newenvironment{reference}{}{}
\newenvironment{slogan}{\comment}{\endcomment}
\newenvironment{history}{\comment}{\endcomment}

% For commutative diagrams we use Xy-pic
\usepackage[all]{xy}

% We use 2cell for 2-commutative diagrams.
\xyoption{2cell}
\UseAllTwocells

% We use multicol for the list of chapters between chapters
\usepackage{multicol}

% This is generally recommended for better output
\usepackage{lmodern}
\usepackage[T1]{fontenc}

% For cross-file-references

% Package for hypertext links:
\usepackage{hyperref}

% For any local file, say "hello.tex" you want to link to please

% Theorem environments.
%
\theoremstyle{plain}
\newtheorem{theorem}[subsection]{Theorem}
\newtheorem{proposition}[subsection]{Proposition}
\newtheorem{lemma}[subsection]{Lemma}

\theoremstyle{definition}
\newtheorem{definition}[subsection]{Definition}
\newtheorem{example}[subsection]{Example}
\newtheorem{exercise}[subsection]{Exercise}
\newtheorem{situation}[subsection]{Situation}

\theoremstyle{remark}
\newtheorem{remark}[subsection]{Remark}
\newtheorem{remarks}[subsection]{Remarks}

\numberwithin{equation}{subsection}

% Macros
%
\def\lim{\mathop{\mathrm{lim}}\nolimits}
\def\colim{\mathop{\mathrm{colim}}\nolimits}
\def\Spec{\mathop{\mathrm{Spec}}}
\def\Hom{\mathop{\mathrm{Hom}}\nolimits}
\def\Ext{\mathop{\mathrm{Ext}}\nolimits}
\def\SheafHom{\mathop{\mathcal{H}\!\mathit{om}}\nolimits}
\def\SheafExt{\mathop{\mathcal{E}\!\mathit{xt}}\nolimits}
\def\Sch{\mathit{Sch}}
\def\Mor{\mathop{\mathrm{Mor}}\nolimits}
\def\Ob{\mathop{\mathrm{Ob}}\nolimits}
\def\Sh{\mathop{\mathit{Sh}}\nolimits}
\def\NL{\mathop{N\!L}\nolimits}
\def\CH{\mathop{\mathrm{CH}}\nolimits}
\def\proetale{{pro\text{-}\acute{e}tale}}
\def\etale{{\acute{e}tale}}
\def\QCoh{\mathit{QCoh}}
\def\Ker{\mathop{\mathrm{Ker}}}
\def\Im{\mathop{\mathrm{Im}}}
\def\Coker{\mathop{\mathrm{Coker}}}
\def\Coim{\mathop{\mathrm{Coim}}}

% Boxtimes
%
\DeclareMathSymbol{\boxtimes}{\mathbin}{AMSa}{"02}

%
% Macros for moduli stacks/spaces
%
\def\QCohstack{\mathcal{QC}\!\mathit{oh}}
\def\Cohstack{\mathcal{C}\!\mathit{oh}}
\def\Spacesstack{\mathcal{S}\!\mathit{paces}}
\def\Quotfunctor{\mathrm{Quot}}
\def\Hilbfunctor{\mathrm{Hilb}}
\def\Curvesstack{\mathcal{C}\!\mathit{urves}}
\def\Polarizedstack{\mathcal{P}\!\mathit{olarized}}
\def\Complexesstack{\mathcal{C}\!\mathit{omplexes}}
% \Pic is the operator that assigns to X its picard group, usage \Pic(X)
% \Picardstack_{X/B} denotes the Picard stack of X over B
% \Picardfunctor_{X/B} denotes the Picard functor of X over B
\def\Pic{\mathop{\mathrm{Pic}}\nolimits}
\def\Picardstack{\mathcal{P}\!\mathit{ic}}
\def\Picardfunctor{\mathrm{Pic}}
\def\Deformationcategory{\mathcal{D}\!\mathit{ef}}

\setlength{\emergencystretch}{3em}
\begin{document}
\title[Stacks Project, Tag 00OC]{1626 --- Auslander\textendash{}Buchsbaum\textendash{}Serre classification of Noetherian local rings with finite residue-field projective dimension and exact global dimension}
\maketitle
\noindent Copyright (C) 2005--2025 Johan de Jong. Stacks Project authors. Source: \href{https://stacks.math.columbia.edu/tag/00OC}{Tag 00OC}.

\noindent Editorial difficulty note (not author text). Difficulty H2: The classification links finite residue-field projective dimension to regularity using a minimal-resolution/Koszul comparison injectivity argument and determinantal depth bounds, then obtains the exact global dimension for arbitrary modules through MCM syzygies and transfinite filtrations. This coupled structural homological argument is substantially beyond ordinary finite-module resolution calculations..

\noindent Editorial provenance note (not author text). History: excerpt from revision \texttt{a0444\allowbreak 6e57e\allowbreak c1fbc\allowbreak 252a8\allowbreak 71afc\allowbreak ec775\allowbreak 2fb28\allowbreak 07b14}, algebra.tex, lines 27095--27123. Reference presentation was adapted; original-author context and core support blocks were retained or added. The mathematical wording is unchanged. Licensed under GNU FDL 1.2 or later; no invariant sections or cover texts. See ../licenses/COPYING. Context blocks reproduce author definitions and supporting results; they are not additional counted problems.

\begin{situation}
\label{situation-complex}
Here $R$ is a ring, and we have a complex
$$
0
\to
R^{n_e}
\xrightarrow{\varphi_e}
R^{n_{e-1}}
\xrightarrow{\varphi_{e-1}}
\ldots
\xrightarrow{\varphi_{i + 1}}
R^{n_i}
\xrightarrow{\varphi_i}
R^{n_{i-1}}
\xrightarrow{\varphi_{i-1}}
\ldots
\xrightarrow{\varphi_1}
R^{n_0}
$$
In other words we require $\varphi_i \circ \varphi_{i + 1} = 0$
for $i = 1, \ldots, e - 1$.
\end{situation}

\begin{definition}
\label{definition-rank}
Let $R$ be a ring. Suppose that $\varphi : R^m \to R^n$ is a map
of finite free modules.
\begin{enumerate}
\item The {\it rank} of $\varphi$ is the maximal $r$ such that
$\wedge^r \varphi : \wedge^r R^m \to \wedge^r R^n$ is nonzero.
\item We let $I(\varphi) \subset R$ be the ideal generated by
the $r \times r$ minors of the matrix of $\varphi$, where $r$
is the rank as defined above.
\end{enumerate}
\end{definition}

\begin{definition}
\label{definition-regular-local}
Let $(R, \mathfrak m)$ be a Noetherian local ring of dimension $d$.
\begin{enumerate}
\item A {\it system of parameters of $R$} is a sequence of elements
$x_1, \ldots, x_d \in \mathfrak m$ which generates an ideal of
definition of $R$,
\item if there exist $x_1, \ldots, x_d \in \mathfrak m$
such that $\mathfrak m = (x_1, \ldots, x_d)$ then we call
$R$ a {\it regular local ring} and $x_1, \ldots, x_d$ a {\it regular
system of parameters}.
\end{enumerate}
\end{definition}

\begin{lemma}
\label{lemma-dimension-shift}
Let $R$ be a Noetherian local Cohen-Macaulay ring of dimension $d$.
Let $0 \to K \to R^{\oplus n} \to M \to 0$ be an exact sequence of
$R$-modules. Then either $M = 0$, or $\text{depth}(K) > \text{depth}(M)$, or
$\text{depth}(K) = \text{depth}(M) = d$.
\end{lemma}

\begin{definition}
\label{definition-CM}
Let $R$ be a Noetherian local ring.
Let $M$ be a finite $R$-module.
We say $M$ is {\it Cohen-Macaulay}
if $\dim(\text{Supp}(M)) = \text{depth}(M)$.
\end{definition}

\begin{definition}
\label{definition-maximal-CM}
Let $R$ be a Noetherian local ring.
A finite module $M$ over $R$ is called a {\it maximal Cohen-Macaulay}
module if $\text{depth}(M) = \dim(R)$.
\end{definition}

\begin{definition}
\label{definition-finite-proj-dim}
Let $R$ be a ring. Let $M$ be an $R$-module. We say $M$ has
{\it finite projective dimension} if it has a finite length
resolution by projective $R$-modules. The minimal length of such a
resolution is called the {\it projective dimension}
of $M$.
\end{definition}

\begin{definition}
\label{definition-regular-sequence}
Let $R$ be a ring. Let $M$ be an $R$-module. A sequence of elements
$f_1, \ldots, f_r$ of $R$ is called an {\it $M$-regular sequence}
if the following conditions hold:
\begin{enumerate}
\item $f_i$ is a nonzerodivisor on
$M/(f_1, \ldots, f_{i - 1})M$
for each $i = 1, \ldots, r$, and
\item the module $M/(f_1, \ldots, f_r)M$ is not zero.
\end{enumerate}
If $I$ is an ideal of $R$ and $f_1, \ldots, f_r \in I$
then we call $f_1, \ldots, f_r$ an {\it $M$-regular sequence
in $I$}. If $M = R$, we call $f_1, \ldots, f_r$ simply a
{\it regular sequence} (in $I$).
\end{definition}

\begin{definition}
\label{definition-depth}
Let $R$ be a ring, and $I \subset R$ an ideal. Let $M$ be a finite $R$-module.
The {\it $I$-depth} of $M$, denoted $\text{depth}_I(M)$, is defined as follows:
\begin{enumerate}
\item if $IM \not = M$, then $\text{depth}_I(M)$ is the supremum in
$\{0, 1, 2, \ldots, \infty\}$ of the lengths of $M$-regular sequences in $I$,
\item if $IM = M$ we set $\text{depth}_I(M) = \infty$.
\end{enumerate}
If $(R, \mathfrak m)$ is local we call $\text{depth}_{\mathfrak m}(M)$ simply
the {\it depth} of $M$.
\end{definition}

\begin{definition}
\label{definition-finite-gl-dim}
Let $R$ be a ring. The ring
$R$ is said to have {\it finite global dimension}
if there exists an integer $n$ such that
every $R$-module has a resolution by
projective $R$-modules of length at most $n$.
The minimal such $n$ is then called the {\it global dimension}
of $R$.
\end{definition}

\begin{lemma}
\label{lemma-length-resolution-residue-field}
Suppose that $R$ is a Noetherian local ring
with maximal ideal $\mathfrak m$ and
residue field $\kappa$. In this case
the projective dimension of $\kappa$ is
$\geq \dim_\kappa \mathfrak m / \mathfrak m^2$.
\end{lemma}

\begin{proof}
Let $x_1 , \ldots, x_n$ be elements of $\mathfrak m$
whose images in $\mathfrak m / \mathfrak m^2$ form a basis.
Consider the {\it Koszul complex} on $x_1, \ldots, x_n$.
This is the complex
$$
0 \to \wedge^n R^n \to \wedge^{n-1} R^n \to \wedge^{n-2} R^n \to
\ldots \to \wedge^i R^n \to \ldots \to R^n \to R
$$
with maps given by
$$
e_{j_1} \wedge \ldots \wedge e_{j_i}
\longmapsto
\sum_{a = 1}^i (-1)^{a + 1} x_{j_a} e_{j_1} \wedge \ldots
\wedge \hat e_{j_a} \wedge \ldots \wedge e_{j_i}
$$
It is easy to see that this is a complex $K_{\bullet}(R, x_{\bullet})$.
Note that the cokernel of the last map of $K_{\bullet}(R, x_{\bullet})$
is $\kappa$ by Lemma \href{https://stacks.math.columbia.edu/tag/00DV}{lemma NAK (Tag 00DV)} part (8).

\medskip\noindent
If $\kappa$ has finite projective dimension $d$, then we can find
a resolution $F_{\bullet} \to \kappa$ by finite free $R$-modules
of length $d$
(Lemma \href{https://stacks.math.columbia.edu/tag/0CXF}{lemma what kind of resolutions Noetherian local (Tag 0CXF)}).
By Lemma \ref{lemma-add-trivial-complex}
we may assume all the maps in the complex $F_{\bullet}$
have the property that $\Im(F_i \to F_{i-1})
\subset \mathfrak m F_{i-1}$, because removing a trivial
summand from the resolution can at worst shorten the resolution.
By Lemma \href{https://stacks.math.columbia.edu/tag/00LS}{lemma compare resolutions (Tag 00LS)} we can find a map
of complexes $\alpha : K_{\bullet}(R, x_{\bullet}) \to F_{\bullet}$
inducing the identity on $\kappa$. We will prove by induction
that the maps $\alpha_i : \wedge^i R^n = K_i(R, x_{\bullet}) \to F_i$
have the property that
$\alpha_i \otimes \kappa : \wedge^i \kappa^n \to F_i \otimes \kappa$
are injective. This shows that $F_n \not = 0$ and hence $d \geq n$
as desired.

\medskip\noindent
The result is clear for $i = 0$ because the composition
$R \xrightarrow{\alpha_0} F_0 \to \kappa$ is nonzero.
Note that $F_0$ must have rank $1$ since
otherwise the map $F_1 \to F_0$ whose cokernel is a single
copy of $\kappa$ cannot have image contained in $\mathfrak m F_0$.

\medskip\noindent
Next we check the case $i = 1$ as we feel that it is instructive;
the reader can skip this as the induction step will deduce the $i = 1$
case from the case $i = 0$. We saw above that
$F_0 = R$ and $F_1 \to F_0 = R$ has image $\mathfrak m$.
We have a commutative diagram
$$
\begin{matrix}
R^n & = & K_1(R, x_{\bullet}) & \to & K_0(R, x_{\bullet}) & = & R \\
& & \downarrow & & \downarrow & & \downarrow \\
& & F_1 & \to & F_0 & = & R
\end{matrix}
$$
where the rightmost vertical arrow is given by multiplication
by a unit. Hence we see that the image of the composition
$R^n \to F_1 \to F_0 = R$ is also equal to $\mathfrak m$.
Thus the map $R^n \otimes \kappa \to F_1 \otimes \kappa$
has to be injective since $\dim_\kappa (\mathfrak m / \mathfrak m^2) = n$.

\medskip\noindent
Let $i \geq 1$ and assume injectivity of $\alpha_j \otimes \kappa$ has been
proved for all $j \leq i - 1$. Consider the commutative diagram
$$
\begin{matrix}
\wedge^i R^n & = & K_i(R, x_{\bullet}) & \to & K_{i-1}(R, x_{\bullet})
& = & \wedge^{i-1} R^n \\
& & \downarrow & & \downarrow & & \\
& & F_i & \to & F_{i-1} & &
\end{matrix}
$$
We know that $\wedge^{i-1} \kappa^n \to F_{i-1} \otimes \kappa$
is injective. This proves that
$\wedge^{i-1} \kappa^n \otimes_{\kappa} \mathfrak m/\mathfrak m^2
\to F_{i-1} \otimes \mathfrak m/\mathfrak m^2$ is injective.
Also, by our choice of the complex, $F_i$ maps into
$\mathfrak mF_{i-1}$, and similarly for the Koszul complex.
Hence we get a commutative diagram
$$
\begin{matrix}
\wedge^i \kappa^n & \to &
\wedge^{i-1} \kappa^n \otimes \mathfrak m/\mathfrak m^2 \\
\downarrow & & \downarrow \\
F_i \otimes \kappa & \to & F_{i-1} \otimes \mathfrak m/\mathfrak m^2
\end{matrix}
$$
At this point it suffices to verify the map
$\wedge^i \kappa^n \to
\wedge^{i-1} \kappa^n \otimes \mathfrak m/\mathfrak m^2$
is injective, which can be done by hand.
\end{proof}


\begin{lemma}
\label{lemma-dim-gl-dim}
Let $R$ be a Noetherian local ring.
Suppose that the residue field $\kappa$ has finite
projective dimension $n$ over $R$.
In this case $\dim(R) \geq n$.
\end{lemma}

\begin{proof}
Let $F_{\bullet}$ be a finite resolution of $\kappa$ by finite free
$R$-modules (Lemma \href{https://stacks.math.columbia.edu/tag/0CXF}{lemma what kind of resolutions Noetherian local (Tag 0CXF)}).
By Lemma \ref{lemma-add-trivial-complex}
we may assume all the maps in the complex $F_{\bullet}$
have to property that $\Im(F_i \to F_{i-1})
\subset \mathfrak m F_{i-1}$, because removing a trivial
summand from the resolution can at worst shorten the resolution.
Say $F_n \not = 0$ and $F_i = 0$ for $i > n$, so that
the projective dimension of $\kappa$ is $n$.
By Proposition \ref{proposition-what-exact} we see that
$\text{depth}_{I(\varphi_n)}(R) \geq n$ since $I(\varphi_n)$
cannot equal $R$ by our choice of the complex.
Thus by Lemma \ref{lemma-bound-depth} also $\dim(R) \geq n$.
\end{proof}


\begin{proposition}
\label{proposition-regular-finite-gl-dim}
Let $R$ be a regular local ring of dimension $d$.
Every finite $R$-module $M$ of depth $e$ has a finite free
resolution
$$
0 \to F_{d-e} \to \ldots \to F_0 \to M \to 0.
$$
In particular a regular local ring has global dimension $\leq d$.
\end{proposition}

\begin{proof}
The first part holds in view of Lemma \ref{lemma-regular-mcm-free}
and Lemma \ref{lemma-mcm-resolution}. The last part follows from this
and Lemma \ref{lemma-finite-gl-dim}.
\end{proof}


\begin{lemma}
\label{lemma-regular-mcm-free}
Let $R$ be a regular local ring.
Any maximal Cohen-Macaulay module over $R$ is free.
\end{lemma}

\begin{proof}
Let $M$ be a maximal Cohen-Macaulay module over $R$.
Let $x \in \mathfrak m$ be part of a regular sequence
generating $\mathfrak m$. Then $x$ is a nonzerodivisor
on $M$ by Proposition \ref{proposition-CM-module}, and
$M/xM$ is a maximal Cohen-Macaulay module over $R/xR$.
By induction on $\dim(R)$ we see that $M/xM$ is free.
We win by Lemma \ref{lemma-free-mod-x}.
\end{proof}


\begin{lemma}
\label{lemma-mcm-resolution}
Let $R$ be a local Noetherian Cohen-Macaulay ring of dimension $d$.
Let $M$ be a finite $R$-module of depth $e$.
There exists an exact complex
$$
0 \to K \to F_{d-e-1} \to \ldots \to F_0 \to M \to 0
$$
with each $F_i$ finite free and $K$ maximal Cohen-Macaulay.
\end{lemma}

\begin{proof}
Immediate from the definition and Lemma \ref{lemma-dimension-shift}.
\end{proof}


\begin{lemma}
\label{lemma-finite-gl-dim}
Let $R$ be a ring. The following are equivalent
\begin{enumerate}
\item $R$ has finite global dimension $\leq n$,
\item every finite $R$-module has projective dimension $\leq n$, and
\item every cyclic $R$-module $R/I$ has projective dimension $\leq n$.
\end{enumerate}
\end{lemma}

\begin{proof}
It is clear that (1) $\Rightarrow$ (2) and (2) $\Rightarrow$ (3).
Assume (3). Choose a set $E \subset M$ of generators of $M$.
Choose a well ordering on $E$. For $e \in E$ denote
$M_e$ the submodule of $M$ generated by the elements $e' \in E$
with $e' \leq e$. Then $M = \bigcup_{e \in E} M_e$.
Note that for each $e \in E$ the quotient
$$
M_e/\bigcup\nolimits_{e' < e} M_{e'}
$$
is either zero or generated by one element, hence has projective
dimension $\leq n$ by (3). By Lemma \ref{lemma-colimit-projective-dimension}
this means that $M$ has projective dimension $\leq n$.
\end{proof}


\begin{lemma}
\label{lemma-regular-ring-CM}
Let $R$ be a regular local ring and let
$x_1, \ldots, x_d$ be a minimal set of generators
for the maximal ideal $\mathfrak m$. Then
$x_1, \ldots, x_d$ is a regular sequence, and
each $R/(x_1, \ldots, x_c)$ is a regular local ring
of dimension $d - c$. In particular $R$ is Cohen-Macaulay.
\end{lemma}

\begin{proof}
Note that $R/x_1R$ is a Noetherian local ring of dimension $\geq d - 1$
by Lemma \href{https://stacks.math.columbia.edu/tag/00KW}{lemma one equation (Tag 00KW)} with $x_2, \ldots, x_d$
generating the maximal ideal. Hence it is a regular local ring by definition.
Since $R$ is a domain by Lemma \ref{lemma-regular-domain}
$x_1$ is a nonzerodivisor.
\end{proof}


\begin{lemma}
\label{lemma-regular-domain}
Any regular local ring is a domain.
\end{lemma}

\begin{proof}
We will use that $\bigcap \mathfrak m^n = 0$
by Lemma \href{https://stacks.math.columbia.edu/tag/00IP}{lemma intersect powers ideal module zero (Tag 00IP)}.
Let $f, g \in R$ such that $fg = 0$.
Suppose that $f \in \mathfrak m^a$ and
$g \in \mathfrak m^b$, with $a, b$ maximal.
Since $fg = 0 \in \mathfrak m^{a + b + 1}$
we see from the result of Lemma \ref{lemma-regular-graded}
that either $f \in \mathfrak m^{a + 1}$ or
$g \in \mathfrak m^{b + 1}$. Contradiction.
\end{proof}


\begin{lemma}
\label{lemma-regular-graded}
Let $(R, \mathfrak m, \kappa)$ be a regular local ring of dimension $d$.
The graded ring $\bigoplus \mathfrak m^n / \mathfrak m^{n + 1}$
is isomorphic to the graded polynomial algebra
$\kappa[X_1, \ldots, X_d]$.
\end{lemma}

\begin{proof}
Let $x_1, \ldots, x_d$ be a minimal set of generators
for the maximal ideal $\mathfrak m$, see
Definition \ref{definition-regular-local}.
There is a surjection $\kappa[X_1, \ldots, X_d]
\to \bigoplus \mathfrak m^n/\mathfrak m^{n + 1}$,
which maps $X_i$ to the class of $x_i$ in $\mathfrak m/\mathfrak m^2$.
Since $d(R) = d$ by
Proposition \href{https://stacks.math.columbia.edu/tag/00KQ}{proposition dimension (Tag 00KQ)}
we know that the numerical
polynomial $n \mapsto \dim_\kappa \mathfrak m^n/\mathfrak m^{n + 1}$
has degree $d - 1$. By Lemma \href{https://stacks.math.columbia.edu/tag/00K3}{lemma quotient smaller d (Tag 00K3)} we
conclude that the surjection $\kappa[X_1, \ldots, X_d]
\to \bigoplus \mathfrak m^n/\mathfrak m^{n + 1}$ is an isomorphism.
\end{proof}


\begin{lemma}
\label{lemma-free-mod-x}
Let $R$ be a Noetherian local ring.
Let $x \in \mathfrak m$.
Let $M$ be a finite $R$-module such that
$x$ is a nonzerodivisor on $M$ and
$M/xM$ is free over $R/xR$.
Then $M$ is free over $R$.
\end{lemma}

\begin{proof}
Let $m_1, \ldots, m_r$ be elements of $M$ which map to
a $R/xR$-basis of $M/xM$. By Nakayama's Lemma \href{https://stacks.math.columbia.edu/tag/00DV}{lemma NAK (Tag 00DV)}
$m_1, \ldots, m_r$ generate $M$. If $\sum a_i m_i = 0$
is a relation, then $a_i \in xR$ for all $i$. Hence
$a_i = b_i x$ for some $b_i \in R$. Hence
the kernel $K$ of $R^r \to M$ satisfies $xK = K$
and hence is zero by Nakayama's lemma.
\end{proof}


\begin{proposition}
\label{proposition-CM-module}
Let $R$ be a Noetherian local ring, with maximal ideal $\mathfrak m$.
Let $M$ be a Cohen-Macaulay module over $R$ whose support has dimension $d$.
Suppose that $g_1, \ldots, g_c$ are elements of
$\mathfrak m$ such that $\dim(\text{Supp}(M/(g_1, \ldots, g_c)M))
= d - c$. Then $g_1, \ldots, g_c$ is an $M$-regular sequence,
and can be extended to a maximal $M$-regular sequence.
\end{proposition}

\begin{proof}
Let $Z = \text{Supp}(M) \subset \Spec(R)$.
By Lemma \href{https://stacks.math.columbia.edu/tag/00KW}{lemma one equation (Tag 00KW)} in the chain
$Z \supset Z \cap V(g_1) \supset \ldots \supset Z \cap V(g_1, \ldots, g_c)$
each step decreases the dimension at most by $1$. Hence by assumption
each step decreases the dimension by exactly $1$ each time. Thus we
may successively apply Lemma \href{https://stacks.math.columbia.edu/tag/00N5}{lemma CM one g (Tag 00N5)} to the modules
$M/(g_1, \ldots, g_i)$ and the element $g_{i + 1}$.

\medskip\noindent
To extend $g_1, \ldots, g_c$ by one element if $c < d$ we simply
choose an element $g_{c + 1} \in \mathfrak m$ which is not
in any of the finitely many minimal primes of $Z \cap V(g_1, \ldots, g_c)$,
using Lemma \href{https://stacks.math.columbia.edu/tag/00DS}{lemma silly (Tag 00DS)}.
\end{proof}


\begin{lemma}
\label{lemma-bound-depth}
Let $(R, \mathfrak m)$ be a Noetherian local ring.
Let $M$ be a nonzero finite $R$-module.
Then $\dim(\text{Supp}(M)) \geq \text{depth}(M)$.
\end{lemma}

\begin{proof}
The proof is by induction on $\dim(\text{Supp}(M))$.
If $\dim(\text{Supp}(M)) = 0$, then
$\text{Supp}(M) = \{\mathfrak m\}$, whence $\text{Ass}(M) = \{\mathfrak m\}$
(by Lemmas \href{https://stacks.math.columbia.edu/tag/0586}{lemma ass support (Tag 0586)} and \href{https://stacks.math.columbia.edu/tag/0587}{lemma ass zero (Tag 0587)}), and hence
the depth of $M$ is zero for example by
Lemma \href{https://stacks.math.columbia.edu/tag/00LL}{lemma ideal nonzerodivisor (Tag 00LL)}.
For the induction step we assume $\dim(\text{Supp}(M)) > 0$.
Let $f_1, \ldots, f_d$ be a sequence of elements of $\mathfrak m$
such that $f_i$ is a nonzerodivisor on $M/(f_1, \ldots, f_{i - 1})M$.
According to Lemma \href{https://stacks.math.columbia.edu/tag/0AUI}{lemma depth weak sequence (Tag 0AUI)} it suffices to prove
$\dim(\text{Supp}(M)) \geq d$. We may assume
$d > 0$ otherwise the lemma holds. By
Lemma \href{https://stacks.math.columbia.edu/tag/0B52}{lemma one equation module (Tag 0B52)}
we have $\dim(\text{Supp}(M/f_1M)) = \dim(\text{Supp}(M)) - 1$.
By induction we conclude $\dim(\text{Supp}(M/f_1M)) \geq d - 1$
as desired.
\end{proof}


\begin{lemma}
\label{lemma-depth-in-ses}
Let $R$ be a local Noetherian ring. Let $0 \to N' \to N \to N'' \to 0$
be a short exact sequence of nonzero finite $R$-modules.
\begin{enumerate}
\item
$\text{depth}(N) \geq \min\{\text{depth}(N'), \text{depth}(N'')\}$
\item
$\text{depth}(N'') \geq \min\{\text{depth}(N), \text{depth}(N') - 1\}$
\item
$\text{depth}(N') \geq \min\{\text{depth}(N), \text{depth}(N'') + 1\}$
\end{enumerate}
\end{lemma}

\begin{proof}
Use the characterization of depth using the Ext groups
$\Ext^i(\kappa, N)$, see Lemma \href{https://stacks.math.columbia.edu/tag/00LW}{lemma depth ext (Tag 00LW)},
and use the long exact cohomology sequence
$$
\begin{matrix}
0
\to \Hom_R(\kappa, N')
\to \Hom_R(\kappa, N)
\to \Hom_R(\kappa, N'')
\\
\phantom{0\ }
\to \Ext^1_R(\kappa, N')
\to \Ext^1_R(\kappa, N)
\to \Ext^1_R(\kappa, N'')
\to \ldots
\end{matrix}
$$
from Lemma \href{https://stacks.math.columbia.edu/tag/00LU}{lemma long exact seq ext (Tag 00LU)}.
\end{proof}


\begin{proposition}
\label{proposition-what-exact}
\begin{reference}
\href{https://stacks.math.columbia.edu/bibliography}{Bibliography: WhatExact (Corollary 1)}
\end{reference}
In Situation \ref{situation-complex}, suppose $R$ is
a local Noetherian ring. The following are equivalent
\begin{enumerate}
\item $0 \to R^{n_e} \to R^{n_{e-1}} \to \ldots \to R^{n_0}$
is exact at $R^{n_e}, \ldots, R^{n_1}$, and
\item for all $i$, $1 \leq i \leq e$
the following two conditions are satisfied:
\begin{enumerate}
\item $\text{rank}(\varphi_i) = r_i$ where
$r_i = n_i - n_{i + 1} + \ldots + (-1)^{e-i-1} n_{e-1} + (-1)^{e-i} n_e$,
\item $I(\varphi_i) = R$, or $I(\varphi_i)$ contains a
regular sequence of length $i$.
\end{enumerate}
\end{enumerate}
\end{proposition}

\begin{proof}
If for some $i$ some matrix coefficient of $\varphi_i$
is not in $\mathfrak m$, then we apply Lemma \ref{lemma-add-trivial-complex}.
It is easy to see that the proposition for a complex and
for the same complex with a trivial complex added to it
are equivalent. Thus we may assume that all matrix entries
of each $\varphi_i$ are elements of the maximal ideal.
We may also assume that $e \geq 1$.

\medskip\noindent
Assume the complex is exact at $R^{n_e}, \ldots, R^{n_1}$.
Let $\mathfrak q \in \text{Ass}(R)$.
Note that the ring $R_{\mathfrak q}$ has depth $0$
and that the complex remains exact after localization at $\mathfrak q$.
We apply Lemmas \ref{lemma-exact-depth-zero-local} and
\ref{lemma-trivial-case-exact} to the localized complex
over $R_{\mathfrak q}$. We conclude that
$\varphi_{i, \mathfrak q}$ has rank $r_i$ for all $i$.
Since $R \to \bigoplus_{\mathfrak q \in \text{Ass}(R)} R_\mathfrak q$
is injective (Lemma \href{https://stacks.math.columbia.edu/tag/0311}{lemma zero at ass zero (Tag 0311)}), we conclude that
$\varphi_i$ has rank $r_i$ over $R$ by the definition of rank as given
in Definition \ref{definition-rank}. Therefore we see that
$I(\varphi_i)_\mathfrak q = I(\varphi_{i, \mathfrak q})$
as the ranks do not change. Since all of the ideals
$I(\varphi_i)_{\mathfrak q}$, $e \geq i \geq 1$
are equal to $R_{\mathfrak q}$ (by the lemmas referenced above)
we conclude none of the ideals $I(\varphi_i)$ is contained in $\mathfrak q$.
This implies that $I(\varphi_e)I(\varphi_{e-1})\ldots I(\varphi_1)$
is not contained in any of the associated primes
of $R$. By Lemma \href{https://stacks.math.columbia.edu/tag/00DS}{lemma silly (Tag 00DS)} we may choose
$x \in I(\varphi_e)I(\varphi_{e - 1})\ldots I(\varphi_1)$,
$x \not \in \mathfrak q$ for all $\mathfrak q \in \text{Ass}(R)$.
Observe that $x$ is a nonzerodivisor (Lemma \href{https://stacks.math.columbia.edu/tag/00LD}{lemma ass zero divisors (Tag 00LD)}).
According to Lemma \ref{lemma-div-x-exact-one-less}
the complex $0 \to (R/xR)^{n_e} \to \ldots \to (R/xR)^{n_1}$ is exact
at $(R/xR)^{n_e}, \ldots, (R/xR)^{n_2}$. By induction
on $e$ all the ideals $I(\varphi_i)/xR$ have a regular
sequence of length $i - 1$. This proves that $I(\varphi_i)$
contains a regular sequence of length $i$.

\medskip\noindent
Assume (2)(a) and (2)(b) hold. We will prove that (1) holds by
induction on $\dim(R)$. If $\dim(R) = 0$, then we must have
$I(\varphi_i) = R$ for $1 \leq i \leq e$ by (2)(b). Since the
coefficients of $\varphi_i$ are contained in the maximal ideal
this can happen only if $r_i = 0$ for all $i$. By (2)(a) we
conclude that $e = 0$ and (1) holds. Assume $\dim(R) > 0$.
We claim that for any prime $\mathfrak p \subset R$ conditions
(2)(a) and (2)(b) hold for the complex
$0 \to R_\mathfrak p^{n_e} \to R_\mathfrak p^{n_{e - 1}} \to \ldots \to
R_\mathfrak p^{n_0}$ with maps $\varphi_{i, \mathfrak p}$
over $R_\mathfrak p$. Namely, since $I(\varphi_i)$ contains a
nonzero divisor, the image of $I(\varphi_i)$ in $R_\mathfrak p$
is nonzero. This implies that the rank of $\varphi_{i, \mathfrak p}$
is the same as the rank of $\varphi_i$: the rank as defined above
of a matrix $\varphi$ over a ring $R$ can only drop when passing
to an $R$-algebra $R'$ and this happens if and only if $I(\varphi)$
maps to zero in $R'$. Thus (2)(a) holds. Having said this
we know that $I(\varphi_{i, \mathfrak p}) = I(\varphi_i)_\mathfrak p$
and we see that (2)(b) is preserved under localization as well.
By induction on the dimension of $R$ we may assume the complex
is exact when localized at any nonmaximal prime $\mathfrak p$ of $R$.
Thus $\Ker(\varphi_i)/\Im(\varphi_{i + 1})$ has support contained in
$\{\mathfrak m\}$ and hence if nonzero has depth $0$.
As $I(\varphi_i) \subset \mathfrak m$ for all $i$ because
of what was said in the first paragraph of the proof, we
see that (2)(b) implies $\text{depth}(R) \geq e$.
By Lemma \ref{lemma-acyclic} we see
that the complex is exact at $R^{n_e}, \ldots, R^{n_1}$
concluding the proof.
\end{proof}


\begin{lemma}
\label{lemma-colimit-projective-dimension}
Let $R$ be a ring. Suppose we have a module $M = \bigcup_{e \in E} M_e$
where the $M_e$ are submodules well-ordered by inclusion. Assume the quotients
$M_e/\bigcup\nolimits_{e' < e} M_{e'}$ have projective dimension $\leq n$.
Then $M$ has projective dimension $\leq n$.
\end{lemma}

\begin{proof}
We will prove this by induction on $n$.

\medskip\noindent
Base case: $n = 0$. Then $P_e = M_e/\bigcup_{e' < e} M_{e'}$ is projective.
Thus we may choose a section $P_e \to M_e$ of the projection $M_e \to P_e$.
We claim that the induced map $\psi : \bigoplus_{e \in E} P_e \to M$ is an
isomorphism. Namely, if $x = \sum x_e \in \bigoplus P_e$ is nonzero,
then we let $e_{max}$ be maximal such that $x_{e_{max}}$ is nonzero
and we conclude that $y = \psi(x) = \psi(\sum x_e)$ is nonzero because
$y \in M_{e_{max}}$ has nonzero image $x_{e_{max}}$ in $P_{e_{max}}$.
On the other hand, let $y \in M$. Then $y \in M_e$ for some $e$.
We show that $y \in \Im(\psi)$ by transfinite induction on $e$.
Let $x_e \in P_e$ be the image of $y$. Then
$y - \psi(x_e) \in \bigcup_{e' < e} M_{e'}$.
By induction hypothesis we conclude that $y - \psi(x_e) \in \Im(\psi)$
hence $y \in \Im(\psi)$. Thus the claim is true and
$\psi$ is an isomorphism. We conclude that $M$ is projective as
a direct sum of projectives, see
Lemma \href{https://stacks.math.columbia.edu/tag/065Q}{lemma direct sum projective (Tag 065Q)}.

\medskip\noindent
If $n > 0$, then for $e \in E$ we denote $F_e$ the free $R$-module
on the set of elements of $M_e$. Then we have a system of
short exact sequences
$$
0 \to K_e \to F_e \to M_e \to 0
$$
over the well-ordered set $E$. Note that the transition maps
$F_{e'} \to F_e$ and $K_{e'} \to K_e$ are injective too.
Set $F = \bigcup F_e$ and $K = \bigcup K_e$. Then
$$
0 \to
K_e/\bigcup\nolimits_{e' < e} K_{e'} \to
F_e/\bigcup\nolimits_{e' < e} F_{e'} \to
M_e/\bigcup\nolimits_{e' < e} M_{e'} \to 0
$$
is a short exact sequence of $R$-modules too and
$F_e/\bigcup_{e' < e} F_{e'}$ is the free $R$-module on the
set of elements in $M_e$ which are not contained in $\bigcup_{e' < e} M_{e'}$.
Hence by
Lemma \ref{lemma-exact-sequence-projective-dimension}
we see that the projective dimension of $K_e/\bigcup_{e' < e} K_{e'}$
is at most $n - 1$. By induction we conclude that $K$ has projective
dimension at most $n - 1$. Whence $M$ has projective dimension at most
$n$ and we win.
\end{proof}


\begin{lemma}
\label{lemma-add-trivial-complex}
Suppose $R$ is a ring. Let
$$
\ldots
\xrightarrow{\varphi_{i + 1}}
R^{n_i}
\xrightarrow{\varphi_i}
R^{n_{i-1}}
\xrightarrow{\varphi_{i-1}}
\ldots
$$
be a complex of finite free $R$-modules. Suppose that for some $i$
some matrix coefficient of the map $\varphi_i$ is invertible.
Then the displayed complex is isomorphic to the direct sum of a complex
$$
\ldots \to
R^{n_{i + 2}} \xrightarrow{\varphi_{i + 2}}
R^{n_{i + 1}} \to
R^{n_i - 1} \to
R^{n_{i - 1} - 1} \to
R^{n_{i - 2}} \xrightarrow{\varphi_{i - 2}}
R^{n_{i - 3}} \to
\ldots
$$
and the complex $\ldots \to 0 \to R \to R \to 0 \to \ldots$
where the map $R \to R$ is the identity map.
\end{lemma}

\begin{proof}
The assumption means, after a change of basis of
$R^{n_i}$ and $R^{n_{i-1}}$ that the first basis
vector of $R^{n_i}$ is mapped via $\varphi_i$ to the first basis
vector of $R^{n_{i-1}}$. Let $e_j$ denote the
$j$th basis vector of $R^{n_i}$ and $f_k$ the $k$th
basis vector of $R^{n_{i-1}}$. Write $\varphi_i(e_j)
= \sum a_{jk} f_k$. So $a_{1k} = 0$ unless $k = 1$
and $a_{11} = 1$. Change basis on $R^{n_i}$ again
by setting $e'_j = e_j - a_{j1} e_1$ for $j > 1$.
After this change of coordinates we have $a_{j1} = 0$
for $j > 1$. Note the image
of $R^{n_{i + 1}} \to R^{n_i}$ is contained in the
subspace spanned by $e_j$, $j > 1$. Note also
that $R^{n_{i-1}} \to R^{n_{i-2}}$ has to annihilate
$f_1$ since it is in the image. These conditions
and the shape of the matrix $(a_{jk})$ for $\varphi_i$
imply the lemma.
\end{proof}


\begin{lemma}
\label{lemma-exact-depth-zero-local}
In Situation \ref{situation-complex}. Suppose $R$ is
a local Noetherian ring with maximal ideal $\mathfrak m$.
Assume $\mathfrak m \in \text{Ass}(R)$, in other words
$R$ has depth $0$. Suppose that
$0 \to R^{n_e} \to R^{n_{e-1}} \to \ldots \to R^{n_0}$
is exact at $R^{n_e}, \ldots, R^{n_1}$.
Then the complex is isomorphic to a direct sum of trivial
complexes.
\end{lemma}

\begin{proof}
Pick $x \in R$, $x \not = 0$, with $\mathfrak m x = 0$.
Let $i$ be the biggest index such that $n_i > 0$.
If $i = 0$, then the statement is true. If
$i > 0$ denote $f_1$ the first basis vector of $R^{n_i}$.
Since $xf_1$ is not mapped to zero by
exactness of the complex we deduce that some matrix
coefficient of the map $R^{n_i} \to R^{n_{i - 1}}$
is not in $\mathfrak m$.
Lemma \ref{lemma-add-trivial-complex} then allows
us to decrease $n_e + \ldots + n_1$. Induction finishes the proof.
\end{proof}


\begin{lemma}
\label{lemma-trivial-case-exact}
In Situation \ref{situation-complex}, suppose the complex is
isomorphic to a direct sum of trivial complexes. Then
we have
\begin{enumerate}
\item the maps $\varphi_i$ have rank
$r_i = n_i - n_{i + 1} + \ldots + (-1)^{e-i-1} n_{e-1} + (-1)^{e-i} n_e$,
\item for all $i$, $1 \leq i \leq e - 1$ we have
$\text{rank}(\varphi_{i + 1}) + \text{rank}(\varphi_i) = n_i$,
\item each $I(\varphi_i) = R$.
\end{enumerate}
\end{lemma}

\begin{proof}
We may assume the complex is the direct sum of trivial
complexes. Then for each $i$ we can split the standard basis
elements of $R^{n_i}$ into those that map to a basis element
of $R^{n_{i-1}}$ and those that are mapped to zero (and these
are mapped onto by basis elements of $R^{n_{i + 1}}$ if $i > 0$).
Using descending
induction starting with $i = e$ it is easy to prove that there
are $r_{i + 1}$-basis elements of $R^{n_i}$ which are mapped
to zero and $r_i$ which are mapped to basis elements of
$R^{n_{i-1}}$. From this the result follows.
\end{proof}


\begin{lemma}
\label{lemma-div-x-exact-one-less}
In Situation \ref{situation-complex}. Suppose $R$ is
a local ring with maximal ideal $\mathfrak m$.
Suppose that $0 \to R^{n_e} \to R^{n_{e-1}} \to \ldots \to R^{n_0}$
is exact at $R^{n_e}, \ldots, R^{n_1}$.
Let $x \in \mathfrak m$ be a nonzerodivisor. The complex
$0 \to (R/xR)^{n_e} \to \ldots \to (R/xR)^{n_1}$
is exact at $(R/xR)^{n_e}, \ldots, (R/xR)^{n_2}$.
\end{lemma}

\begin{proof}
Denote $F_\bullet$ the complex with terms $F_i = R^{n_i}$
and differential given by $\varphi_i$. Then we have a short
exact sequence of complexes
$$
0 \to F_\bullet \xrightarrow{x} F_\bullet \to F_\bullet/xF_\bullet \to 0
$$
Applying the snake lemma we get a long exact sequence
$$
H_i(F_\bullet) \xrightarrow{x} H_i(F_\bullet) \to
H_i(F_\bullet/xF_\bullet) \to H_{i - 1}(F_\bullet)
\xrightarrow{x} H_{i - 1}(F_\bullet)
$$
The lemma follows.
\end{proof}


\begin{lemma}[Acyclicity lemma]
\label{lemma-acyclic}
\begin{reference}
\href{https://stacks.math.columbia.edu/bibliography}{Bibliography: Peskine-Szpiro (Lemma 1.8)}
\end{reference}
Let $R$ be a local Noetherian ring.
Let $0 \to M_e \to M_{e-1} \to \ldots \to M_0$
be a complex of finite $R$-modules.
Assume $\text{depth}(M_i) \geq i$.
Let $i$ be the largest index such that the complex is
not exact at $M_i$. If $i > 0$ then
$\Ker(M_i \to M_{i-1})/\Im(M_{i + 1} \to M_i)$
has depth $\geq 1$.
\end{lemma}

\begin{proof}
Let $H = \Ker(M_i \to M_{i-1})/\Im(M_{i + 1} \to M_i)$ be the
cohomology group in question.
We may break the complex into short exact sequences
$0 \to M_e \to M_{e-1} \to K_{e-2} \to 0$,
$0 \to K_j \to M_j \to K_{j-1} \to 0$, for $i + 2 \leq j \leq e-2 $,
$0 \to K_{i + 1} \to M_{i + 1} \to B_i \to 0$,
$0 \to K_i \to M_i \to M_{i-1}$, and
$0 \to B_i \to K_i \to H \to 0$.
We proceed up through these complexes to
prove the statements about depths, repeatedly using
Lemma \ref{lemma-depth-in-ses}.
First of all, since $\text{depth}(M_e) \geq e$,
and $\text{depth}(M_{e-1}) \geq e-1$ we deduce
that $\text{depth}(K_{e-2}) \geq e - 1$. At this point the
sequences $0 \to K_j \to M_j \to K_{j-1} \to 0$ for $i + 2 \leq j \leq e-2 $
imply similarly that $\text{depth}(K_{j-1}) \geq j$ for
$i + 2 \leq j \leq e-2$. The sequence
$0 \to K_{i + 1} \to M_{i + 1} \to B_i \to 0$
then shows that $\text{depth}(B_i) \geq i + 1$. The sequence
$0 \to K_i \to M_i \to M_{i-1}$ shows that $\text{depth}(K_i) \geq 1$
since $M_i$ has depth $\geq i \geq 1$ by assumption.
The sequence $0 \to B_i \to K_i \to H \to 0$ then
implies the result.
\end{proof}


\begin{lemma}
\label{lemma-exact-sequence-projective-dimension}
Let $R$ be a ring. Let $0 \to M' \to M \to M'' \to 0$ be a short
exact sequence of $R$-modules.
\begin{enumerate}
\item If $M$ has projective dimension $\leq n$ and $M''$
has projective dimension $\leq n + 1$, then $M'$ has projective
dimension $\leq n$.
\item If $M'$ and $M''$ have projective dimension
$\leq n$ then $M$ has projective dimension $\leq n$.
\item If $M'$ has projective dimension $\leq n$ and
$M$ has projective dimension $\leq n + 1$ then
$M''$ has projective dimension $\leq n + 1$.
\end{enumerate}
\end{lemma}

\begin{proof}
Combine the characterization of projective dimension in
Lemma \href{https://stacks.math.columbia.edu/tag/065R}{lemma projective dimension ext (Tag 065R)}
with the long exact sequence of ext groups in
Lemma \href{https://stacks.math.columbia.edu/tag/065P}{lemma reverse long exact seq ext (Tag 065P)}.
\end{proof}


Let $(R, \mathfrak m)$ be a Noetherian local ring.
From the above it is clear that $\mathfrak m$ cannot be
generated by fewer than $\dim(R)$ variables.
By Nakayama's Lemma \href{https://stacks.math.columbia.edu/tag/00DV}{lemma NAK (Tag 00DV)} the minimal number
of generators of $\mathfrak m$ equals $\dim_{\kappa(\mathfrak m)}
\mathfrak m/\mathfrak m^2$. Hence we have the following
fundamental inequality
$$
\dim(R) \leq \dim_{\kappa(\mathfrak m)} \mathfrak m/\mathfrak m^2.
$$
It turns out that the rings where equality holds
have a lot of good properties. They are called
regular local rings.



\begin{proposition}
\label{proposition-finite-gl-dim-regular}
Let $(R, \mathfrak m, \kappa)$ be a Noetherian local ring.
The following are equivalent
\begin{enumerate}
\item $\kappa$ has finite projective dimension as an $R$-module,
\item $R$ has finite global dimension,
\item $R$ is a regular local ring.
\end{enumerate}
Moreover, in this case the global dimension of $R$ equals
$\dim(R) = \dim_\kappa(\mathfrak m/\mathfrak m^2)$.
\end{proposition}

\begin{proof}
We have (3) $\Rightarrow$ (2) by
Proposition \ref{proposition-regular-finite-gl-dim}.
The implication (2) $\Rightarrow$ (1) is trivial.
Assume (1). By Lemmas \ref{lemma-length-resolution-residue-field}
and \ref{lemma-dim-gl-dim} we see that
$\dim(R) \geq \dim_\kappa(\mathfrak m /\mathfrak m^2)$.
Thus $R$ is regular, see Definition
\ref{definition-regular-local} and the discussion preceding it.
Assume the equivalent conditions (1) -- (3) hold.
By Proposition \ref{proposition-regular-finite-gl-dim}
the global dimension of $R$ is at most $\dim(R)$ and by
Lemma \ref{lemma-length-resolution-residue-field}
it is at least $\dim_\kappa(\mathfrak m/\mathfrak m^2)$.
Thus the stated equality holds.
\end{proof}
\end{document}
